\documentclass[11pt,a4paper]{article}

% Packages
\usepackage[utf8]{inputenc}
\usepackage[T1]{fontenc}
\usepackage[margin=0.75in]{geometry}
\usepackage{enumitem}
\usepackage{hyperref}
\usepackage{titlesec}
\usepackage{xcolor}

% Formatting
\pagestyle{empty}
\setlength{\parindent}{0pt}
\setlist[itemize]{leftmargin=*, noitemsep, topsep=3pt}

% Section formatting
\titleformat{\section}{\Large\bfseries}{}{0em}{}[\titlerule]
\titlespacing*{\section}{0pt}{12pt}{6pt}

\titleformat{\subsection}[runin]{\bfseries}{}{0em}{}
\titlespacing*{\subsection}{0pt}{8pt}{0pt}

% Colors
\definecolor{linkcolor}{RGB}{0,102,204}
\hypersetup{
    colorlinks=true,
    urlcolor=linkcolor
}

\begin{document}

% Header
\begin{center}
    {\LARGE \textbf{ERIC MAYEUX}}\\
    \vspace{4pt}
    {\large Gestionnaire de produit -- Robotique \& Automatisation | Product Manager}\\
    \vspace{4pt}
    Montréal, QC, Canada\\
    \href{mailto:mayeux.e@gmail.com}{mayeux.e@gmail.com} | 438-884-7645
\end{center}

\vspace{8pt}

\section*{Profil Professionnel}
\noindent
\textbf{Gestionnaire de produit} avec 10+ ans d'expérience en gestion de \textbf{feuille de route produit}, \textbf{priorisation du backlog} et livraison en environnement Agile. Expert en \textbf{gestion des parties prenantes} internes et externes, traduction des besoins techniques en \textbf{exigences produit} claires. Maîtrise des outils de \textbf{gestion de portefeuille produit} (Jira, Confluence) et de l'\textbf{intelligence artificielle appliquée aux produits} (agents IA, automatisation). Certifié \textbf{PMP} (en cours) et CSM. Bilinguisme français/anglais.

\section*{Compétences Techniques}

\textbf{Gestion de produit:} Feuille de route produit, priorisation du backlog, cahiers de besoins, User Stories, exigences fonctionnelles, indicateurs de performance produit (adoption, satisfaction)

\textbf{Gestion du portefeuille produit (Product Ownership):} Backlog refinement, sprint planning, démos, rétrospectives, acceptance criteria, QA fonctionnelle

\textbf{Intelligence Artificielle:} Agents IA, Playwright MCP, Crew AI, IA appliquée aux produits, RPA, automatisation industrielle

\textbf{Analyse \& Données:} Indicateurs de performance, tableaux de bord (Jira, Datadog, Splunk), analyse quantitative et qualitative, KPIs produit

\textbf{Parties prenantes:} Gestion des parties prenantes internes/externes, communication technique/non-technique, animation de sprint planning et revues trimestrielles

\textbf{Méthodologies:} Agile, Scrum, SAFe, Kanban, TDD, BDD, spec-driven development

\textbf{Architecture \& Cloud:} AWS Cloud Practitioner, CI/CD, Architectures Distribuées, Big Data

\vspace{6pt}

\section*{Réalisations}
\begin{itemize}
    \item Définition et pilotage de la \textbf{feuille de route produit} qualité à la BNC en contexte SAFe
    \item \textbf{Priorisation du backlog} en collaboration avec les PO : alignement valeur, risques, exigences techniques
    \item Mise en place d'\textbf{indicateurs de performance produit} : KPIs release, DORA, DevEx, qualité
    \item \textbf{Traduction des besoins techniques} en exigences fonctionnelles pour des équipes pluridisciplinaires
    \item Intégration d'\textbf{intelligence artificielle} dans des produits : agents IA, Playwright MCP, Crew AI
    \item Animation de \textbf{démos, sprint planning et revues} pour 2 équipes en parallèle (SAFe Release Train)
\end{itemize}

\section*{Expérience Professionnelle}

\subsection*{Scrum Master -- Gestion de projet | Project Manager} \hfill \textit{Novembre 2025 -- Présent}\\
\textbf{Banque Nationale du Canada (BNC)} -- Montréal, QC\\
Pilote la \textbf{feuille de route produit} qualité et la livraison en contexte SAFe.
\begin{itemize}
    \item Collabore avec les PO pour définir les \textbf{exigences produit} et \textbf{prioriser le backlog}
    \item Anime les cérémonies Agile : \textbf{sprint planning, démos, rétrospectives, raffinement}
    \item Suit et communique les \textbf{indicateurs de performance} (KPIs, DORA, DevEx) aux parties prenantes
    \item Intègre des outils d'\textbf{intelligence artificielle} : agents IA, automatisation, Playwright MCP
    \item Crée des \textbf{tableaux de bord} pour le suivi de la santé produit et des pratiques Agiles
    \item Coordonne les dépendances et \textbf{gère les parties prenantes} internes et externes
\end{itemize}

\subsection*{Intégrateur Sénior Qualité | Senior Quality Integrator} \hfill \textit{Octobre 2023 -- Novembre 2025}\\
\textbf{Banque Nationale du Canada (BNC)} -- Montréal, QC\\
\begin{itemize}
    \item Définit les \textbf{exigences fonctionnelles} des stratégies de test et \textbf{feuilles de route} qualité
    \item \textbf{Priorise et gère le backlog} des initiatives qualité pan-train (dont squads Thaïlande)
    \item Intègre Playwright MCP, Crew AI, POC d'outils : \textbf{IA appliquée aux produits} logiciels
\end{itemize}

\subsection*{Intégrateur Qualité -- SDET | Quality Integrator -- Software Development Engineer in Test} \hfill \textit{Janvier 2022 -- Octobre 2023}\\
\textbf{Banque Nationale du Canada (BNC)} via Groupe Astek -- Montréal, QC\\
\begin{itemize}
    \item Traduction des \textbf{exigences produit} en stratégies de test : TDD, BDD, spec-driven
    \item Stack : Selenium, AppliTools, RestAssured, JMeter, Gatling, K6
\end{itemize}

\subsection*{QA Automaticien | QA Automation Engineer} \hfill \textit{Juin 2021 -- Décembre 2021}\\
\textbf{AODocs} -- Paris, France\\
\begin{itemize}
    \item POC Flutter, automatisation non-régression ; \textbf{cahiers de tests fonctionnels}
\end{itemize}

\subsection*{Automaticien CI/CD | DevOps Automation Engineer} \hfill \textit{Janvier 2020 -- Juin 2021}\\
\textbf{ENGIE DIGITAL} via Groupe Hénix -- Paris, France\\
\begin{itemize}
    \item Jenkins, Docker : création CI/CD ; administration Jira, KPIs, \textbf{pilotage produit}
\end{itemize}

\subsection*{Product Owner} \hfill \textit{Juin 2019 -- Septembre 2019}\\
\textbf{Hénix} -- Montrouge, France\\
\begin{itemize}
    \item \textbf{Gestion du backlog}, recette, priorisation, support clients SaaS (Jira/Squash)
\end{itemize}

\subsection*{Développeur Logiciel | Software Developer} \hfill \textit{Mai 2017 -- Mai 2019}\\
\textbf{MEDIAPOST} via Groupe Hénix -- Paris, France\\
\begin{itemize}
    \item Développement back office, webservices ; \textbf{traduction d'exigences} fonctionnelles en code
\end{itemize}

\subsection*{Consultant QA \& Automatisation | QA Automation Consultant} \hfill \textit{Avril 2016 -- Avril 2017}\\
\textbf{ENGIE IT} via Groupe Hénix -- Saint-Ouen, France\\
\begin{itemize}
    \item Coordination opérationnelle, administration HP ALM/UFT, gestion des incidents
\end{itemize}

\section*{Certifications \& Formations}

\textbf{PMP Certification} (en cours) \hfill \textit{2026}

Project Management Professional -- Project Management Institute (PMI)

\textbf{AWS Cloud Practitioner} \hfill \textit{2026}

Amazon Web Services (AWS)

\textbf{Certified Scrum Master (CSM)} \hfill \textit{2024}

Scrum Alliance

\textbf{Jira ACP-600 -- Project Administration in Jira Server} \hfill \textit{2019}

Atlassian Certified Professional

\textbf{Cursus Automaticien et DevOps} \hfill \textit{2019}\\
\textit{AFCEPF -- Montrouge, France}

\textbf{Cursus Architecture logiciel \& Analyste informaticien - Certifié Master II} \hfill \textit{2017}\\
\textit{AFCEPF-HENIX -- Malakoff, France}

\section*{Formation Académique}

\textbf{Licence en Gestion (Bachelor's Degree in Management)} \hfill \textit{2011}\\
École Supérieure de Gestion (E.S.G) -- Paris, France

\section*{Compétences Linguistiques}

\textbf{Français:} Langue maternelle (Native)\\
\textbf{Anglais:} Niveau Avancé (Advanced / Professional Working Proficiency)

\end{document}
